\documentclass[10pt]{beamer}

\usepackage[utf8]{inputenc}
\usepackage[spanish]{babel}
\usepackage{graphicx}

\mode<presentation>
\usetheme{Madrid}
%\usecolortheme[RGB={111,73,135}]{structure}
\usecolortheme[RGB={128,0,0}]{structure}
%\usecolortheme[RGB={0,96,0}]{structure}
%\usecolortheme[RGB={200,0,200}]{structure}
%\usecolortheme[RGB={0,128,0}]{structure}
%\usecolortheme[RGB={0,0,128}]{structure}
\usefonttheme{serif}
\useinnertheme{rectangles}
\useoutertheme{split}
\setbeamercovered{transparent}


% Definiciones para usar luego
\newtheorem{ejemplo}{Ejemplo}
\newtheorem{definicion}{Definición}

\title{Entendiendo y Optimizando MySQL}
\author{Jesús Espino García}
\date{10 de Noviembre de 2010}
\subject{Entendiendo y Optmizando MySQL}

\institute[GUL UC3M]{
  Grupo de Usuarios de Linux\\
  Universidad Carlos III de Madrid.\\
  \ \\
  \includegraphics[height=1.5cm]{imgs/gul}}

\setcounter{tocdepth}{2}

\AtBeginSubsection[]
{
  \begin{frame}<beamer>{Indice}
    \tableofcontents[sectionstyle=show/shaded,subsectionstyle=show/shaded/hide]
  \end{frame}
}

\begin{document}

  \frame{\maketitle}
  
  \include{mysq_arquitectura}
  \include{mysq_esquemas_e_indices}
  \include{mysq_storage_engines}
  \include{mysq_queries}
  \include{mysq_analizando_el_rendimiento}
  \include{mysq_optimizador}
  \include{mysq_cache}
  \include{mysq_escalando}
  \section*{Para terminar}
\subsection*{Otra información útil}
\begin{frame}{Referencias}
  \begin{itemize}
    \item Manual oficial de mysql (mysql.com).
    \item O'Reilly - High Performance MySQL (Second Edition).
    \item Apress - Pro MySQL.
  \end{itemize}
\end{frame}

\subsection*{Cierre}

\begin{frame}{¿Qué queda en el tintero?}
  \begin{itemize}
    \item Alta disponibilidad
    \item Replicación
    \item Optimizaciones de configuración.
    \item Optimizaciones de SSOO.
    \item Optimizaciones de Hardware.
    \item \dots
  \end{itemize}
\end{frame}

\begin{frame}{Dudas}
  \begin{center}
    \dots
  \end{center}
\end{frame}

\begin{frame}{Fin}
  \begin{center}
    Gracias por venir.
  \end{center}
\end{frame}


\end{document}
